\clearpage
\chapter*{2024 Enero}
\addcontentsline{toc}{chapter}{2024 Enero}

\begin{enunciado}
    \hypertarget{2024 A 1}{}
    \textbf{1. (3,5 puntos)} En el conjunto \(\mathbb{R}\) de los números reales se considera la topología \(T\) engendrada por la subbase \(S = \{[a,b] \mid a,b \in \mathbb{Q}\}\). Se pide:
    \begin{enumerate}[label=\alph*)]
        \item Estudiar si \(T_u \subset T\).
        \item Estudiar si \((\mathbb{R},T)\) es conexo.
    \end{enumerate}
\end{enunciado}

\begin{enunciado}
    \hypertarget{2024 A 2}{}
    \textbf{2. (3,5 puntos)} En \(X = \mathbb{R}^2\) se da la topología \(T\) cuya base es
    \[
        B = \{[a,b) \times [c,d) \mid a,b,c,d \in \mathbb{R}\}.
    \]
    Si \(A = \{(x,y) \in \mathbb{R}^2 \mid x+y=0\}\) y \(C = \{(x,y) \in \mathbb{R}^2 \mid x-y=0\}\), se pide:
    \begin{enumerate}[label=\alph*)]
        \item Determinar la topología \(T_A\) y dar una base de la topología \(T_C\).
        \item Estudiar si \((A,T_A)\) y \((C,T_C)\) son espacios topológicos separables.
    \end{enumerate}
\end{enunciado}

\begin{enunciado}
    \hypertarget{2024 A 3}{}
    \textbf{3. (3 puntos)} Sea \((X,T)\) un espacio topológico. Se pide:
    \begin{enumerate}[label=\alph*)]
        \item Definir punto de aglomeración de una sucesión en \((X,T)\).
        \item Estudiar si en un espacio compacto \((X,T)\) toda sucesión tiene al menos un punto de aglomeración.
    \end{enumerate}
\end{enunciado}